% Zumdahl Chapter 15 test -- 20 MCQ + 1 long FRQ + 1 short FRQ. 60 min.
\documentclass{apchem-exam}

\apcunitword{Chapter}
\apcunit{15}
\apcunittitle{Acid--Base Equilibria}
\apcdoctitle{Chapter Test}
\apcsubtitle{Zumdahl \S15.1--15.5 \textbullet\ 60 minutes \textbullet\ $K$ and p$K_a$ values provided with Section II}

\begin{document}
\apctitleblock

\begin{apcnote}[Scope and format]
Covers \S15.1--15.5. Nothing from \S15.6 (polyprotic titrations) is
calculated; you are responsible only for the \emph{shape} of such a curve,
which is asked qualitatively.

This test carries \textbf{one long and one short free response} rather than
the usual two short ones. A full titration is the canonical AP long-response
topic and cannot be assessed in four points.
\end{apcnote}

\examsection{I}{Multiple Choice}{20 questions \textbullet\ 30 minutes \textbullet\ 1 point each}{\nocalculator}

% ---- common ion ------------------------------------------------------------
\begin{mcq}
  Adding \ce{NaF} to a solution of \ce{HF} causes the pH to
  \choices
    \correct{rise, because \ce{F-} shifts the ionization to the left}
    \choice{fall, because \ce{NaF} is a salt}
    \choice{stay the same, because \ce{NaF} is neutral}
    \choice{rise, because \ce{Na+} is basic}
\end{mcq}
\why{The common ion effect is Le Ch\^atelier applied to the acid
ionization.}
\distractor{D}{\ce{Na+} has no acid--base properties at all}

\begin{mcq}
  Which addition would have \emph{no} effect on the pH of a
  \ce{HC2H3O2} solution?
  \choices
    \choice{\ce{NaC2H3O2}}
    \correct{\ce{NaCl}}
    \choice{\ce{HCl}}
    \choice{\ce{NaOH}}
\end{mcq}
\why{Neither \ce{Na+} nor \ce{Cl-} participates in the acetic acid
equilibrium.}

\begin{mcq}
  Which pair constitutes a buffer?
  \choices
    \choice{\ce{HCl} and \ce{NaCl}}
    \choice{\ce{NaOH} and \ce{NaCl}}
    \correct{\ce{NH3} and \ce{NH4Cl}}
    \choice{\ce{HNO3} and \ce{KNO3}}
\end{mcq}
\why{A buffer needs both members of a \emph{weak} conjugate pair.}
\distractor{A}{\ce{Cl-} is the conjugate base of a strong acid and has no
affinity for protons}

\begin{mcq}
  A buffer can also be prepared by mixing a weak acid with
  \choices
    \correct{less than a stoichiometric amount of strong base}
    \choice{exactly a stoichiometric amount of strong base}
    \choice{an excess of strong base}
    \choice{a strong acid}
\end{mcq}
\why{Partial neutralization leaves both \ce{HA} and \ce{A-} in solution.}
\distractor{B}{that is the equivalence point --- only \ce{A-} remains}

% ---- Henderson-Hasselbalch -------------------------------------------------
\begin{mcq}
  The Henderson--Hasselbalch equation is
  \choices
    \correct{pH $=$ p$K_a + \log([\text{base}]/[\text{acid}])$}
    \choice{pH $=$ p$K_a + \log([\text{acid}]/[\text{base}])$}
    \choice{pH $=$ p$K_a - [\text{base}]/[\text{acid}]$}
    \choice{pH $=$ p$K_a \times \log([\text{base}]/[\text{acid}])$}
\end{mcq}
\why{Base over acid; inverting the ratio flips the sign of the correction.}

\begin{mcq}
  A buffer contains equal concentrations of \ce{HC2H3O2} and
  \ce{C2H3O2^-} (p$K_a$ 4.74). Its pH is
  \choices
    \choice{1.00}
    \correct{4.74}
    \choice{7.00}
    \choice{9.26}
\end{mcq}
\why{Equal amounts make the log term zero.}

\begin{mcq}
  A buffer contains ten times as much \ce{A-} as \ce{HA}. Its pH is
  \choices
    \choice{p$K_a - 1$}
    \correct{p$K_a + 1$}
    \choice{p$K_a + 10$}
    \choice{$10\,\mathrm{p}K_a$}
\end{mcq}
\why{$\log 10 = 1$.}

\begin{mcq}
  A buffer whose pH lies below its p$K_a$ contains
  \choices
    \correct{more acid than conjugate base}
    \choice{more conjugate base than acid}
    \choice{equal amounts of both}
    \choice{no conjugate base at all}
\end{mcq}
\why{A negative log term requires the denominator to be larger.}

\begin{mcq}
  Diluting a buffer with water
  \choices
    \correct{leaves the pH nearly unchanged but lowers the capacity}
    \choice{raises the pH and raises the capacity}
    \choice{lowers the pH and lowers the capacity}
    \choice{leaves both the pH and the capacity unchanged}
\end{mcq}
\why{The ratio survives dilution; the absolute amounts do not.}

% ---- buffer action and capacity --------------------------------------------
\begin{mcq}
  When \ce{NaOH} is added to a \ce{HC2H3O2}/\ce{NaC2H3O2} buffer, the
  species consumed is
  \choices
    \correct{\ce{HC2H3O2}}
    \choice{\ce{C2H3O2^-}}
    \choice{\ce{Na+}}
    \choice{water}
\end{mcq}
\why{Added base reacts with the acidic member of the pair.}

\begin{mcq}
  The correct first step in a ``buffer plus strong base'' problem is
  \choices
    \correct{a stoichiometry calculation taken to completion}
    \choice{an ICE table containing the added \ce{OH-}}
    \choice{substituting the original concentrations into
      Henderson--Hasselbalch}
    \choice{calculating $K_b$ for the conjugate base}
\end{mcq}
\why{Strong base does not sit at equilibrium with a weak acid; it consumes
it.}
\distractor{B}{putting strong base into an ICE table is the most common way
to lose the whole question}

\begin{mcq}
  Buffering capacity is greatest when
  \choices
    \correct{the two components are present in equal, large amounts}
    \choice{the conjugate base greatly exceeds the acid}
    \choice{the acid greatly exceeds the conjugate base}
    \choice{the solution is as dilute as possible}
\end{mcq}
\why{A ratio of 1 is the least sensitive to change, and more material means
more to consume.}

\begin{mcq}
  A buffer is useful over a pH range of approximately
  \choices
    \choice{p$K_a \pm 0.1$}
    \correct{p$K_a \pm 1$}
    \choice{p$K_a \pm 7$}
    \choice{any pH at all}
\end{mcq}
\why{Beyond about a 10:1 ratio the minor component is too scarce to absorb
anything.}

\begin{mcq}
  To prepare a buffer of pH 9.2, the best choice of weak acid has p$K_a$
  \choices
    \choice{3.4}
    \choice{4.7}
    \choice{7.5}
    \correct{9.2}
\end{mcq}
\why{Choose p$K_a$ as close to the target pH as possible, so the ratio is
near 1.}

\begin{mcq}
  A buffer fails once
  \choices
    \correct{one of the two components has been essentially used up}
    \choice{the solution has been diluted}
    \choice{the pH reaches exactly p$K_a$}
    \choice{any strong acid at all has been added}
\end{mcq}
\why{With only one member of the pair left there is nothing to absorb
further addition.}

% ---- titrations ------------------------------------------------------------
\begin{mcq}
  The equivalence point of a titration is defined by
  \choices
    \correct{the stoichiometry of the reaction}
    \choice{the pH reaching 7.00}
    \choice{the indicator changing color}
    \choice{the steepest part of the curve}
\end{mcq}
\why{It is where moles of titrant equal moles of analyte --- nothing about
pH is implied.}
\distractor{C}{that is the \emph{end point}}

\begin{mcq}
  \SI{25.0}{\milli\litre} of \SI{0.10}{\Molar} \ce{HCl} is titrated with
  \SI{0.10}{\Molar} \ce{NaOH}. The pH at the equivalence point is
  \choices
    \choice{below 7}
    \correct{exactly 7.00}
    \choice{above 7}
    \choice{impossible to predict}
\end{mcq}
\why{Only \ce{Na+} and \ce{Cl-} remain, and neither reacts with water.}

\begin{mcq}
  A weak acid titrated with a strong base has an equivalence point that is
  \choices
    \choice{acidic}
    \choice{exactly neutral}
    \correct{basic}
    \choice{dependent on the titrant concentration only}
\end{mcq}
\why{The solution at equivalence is the conjugate base of a weak acid, which
hydrolyses.}

\begin{mcq}
  A weak base titrated with a strong acid has an equivalence point that is
  \choices
    \correct{acidic}
    \choice{exactly neutral}
    \choice{basic}
    \choice{at pH equal to p$K_b$}
\end{mcq}
\why{Only the conjugate acid of the weak base remains.}

\begin{mcq}
  Halfway to the equivalence point of a weak acid titration,
  \choices
    \correct{pH $=$ p$K_a$}
    \choice{pH $= 7.00$}
    \choice{pH $=$ p$K_b$}
    \choice{the buffer has failed}
\end{mcq}
\why{Half the acid has become its conjugate base, so the ratio is 1.}

\examsection{II}{Free Response}{2 questions \textbullet\ 30 minutes}{\calculatorok}

\begin{apcnote}[Data]
$K_a$: \ce{HOCl} $3.5\times10^{-8}$ (p$K_a$ 7.46) \textbullet\
\ce{HC2H3O2} $1.8\times10^{-5}$ (p$K_a$ 4.74).
$K_w = 1.0\times10^{-14}$. \\
Indicator ranges: methyl red 4.4--6.2 \textbullet\ bromthymol blue
6.0--7.6 \textbullet\ phenolphthalein 8.2--10.0 \textbullet\ alizarin yellow
10.1--12.0.
\end{apcnote}

% ---- FRQ 1 (long) ----------------------------------------------------------
\frq{long}{10}{Zumdahl \S15.4--15.5 \textbullet\ a complete weak acid titration}

A \SI{25.00}{\milli\litre} sample of \SI{0.100}{\Molar} hypochlorous acid,
\ce{HOCl}, is titrated with \SI{0.100}{\Molar} \ce{NaOH}.

\begin{parts}
  \item Calculate the volume of \ce{NaOH} required to reach the equivalence
        point. \answerblank[4cm]{\SI{25.00}{\milli\litre}}
  \item Calculate the pH before any base is added. \workspace[2.4cm]
        \keyonly{\begin{apcanswer}$x = \sqrt{(3.5\times10^{-8})(0.100)}
        = 5.9\times10^{-5}$; pH $= \mathbf{4.23}$\end{apcanswer}}
  \item Calculate the pH after \SI{12.50}{\milli\litre} of base has been
        added, and explain why this value could have been written down
        without calculation. \answerlines[3]{This is half the equivalence
        volume, so half the \ce{HOCl} has become \ce{OCl-} and
        $[\ce{HOCl}] = [\ce{OCl^-}]$. The log term vanishes and
        pH $=$ p$K_a$ $= 7.46$.}
  \item Calculate the pH after \SI{20.00}{\milli\litre} of base.
        \workspace[2.4cm]
        \keyonly{\begin{apcanswer}\ce{HOCl} $= 2.500 - 2.000 = 0.500$,
        \ce{OCl-} $= 2.000$ mmol;
        pH $= 7.46 + \log(2.000/0.500) = 7.46 + 0.60 =
        \mathbf{8.06}$\end{apcanswer}}
  \item Calculate the pH at the equivalence point. \workspace[3.2cm]
        \keyonly{\begin{apcanswer}All \SI{2.500}{\milli\mole} is \ce{OCl-}
        in \SI{50.00}{\milli\litre}, so $[\ce{OCl^-}] = 0.0500$.
        $K_b = 1.0\times10^{-14}/3.5\times10^{-8} = 2.9\times10^{-7}$;
        $x = \sqrt{(2.9\times10^{-7})(0.0500)} = 1.2\times10^{-4}$;
        pOH $= 3.92$; pH $= \mathbf{10.08}$\end{apcanswer}}
  \item Select an indicator for this titration from the list provided and
        justify your choice. \answerlines[2]{Phenolphthalein (8.2--10.0) ---
        its range brackets the equivalence pH of 10.08 most closely of the
        available choices, so its end point falls near equivalence. Methyl
        red and bromthymol blue would both change color while acid
        remained.}
\end{parts}

\begin{rubric}
  \rpoint{1}{(a) \SI{25.00}{\milli\litre}}
  \rpoint{2}{(b) 1 pt correct $\sqrt{K_aC_0}$ setup; 1 pt pH $= 4.23$}
  \rpoint{2}{(c) 1 pt pH $= 7.46$; 1 pt reasoning that identifies the
    halfway point AND states $[\ce{HA}] = [\ce{A^-}]$ --- naming the value
    with no reasoning earns 1}
  \rpoint{2}{(d) 1 pt correct mmol bookkeeping; 1 pt pH $= 8.06$}
  \rpoint{2}{(e) 1 pt dilution to \SI{50.00}{\milli\litre} AND $K_b$ from
    $K_w/K_a$; 1 pt pH $= 10.08$ via pOH --- reporting $-\log x$ as the pH
    earns 0 for this point}
  \rpoint{1}{(f) phenolphthalein with a justification tied to the
    equivalence pH --- naming it alone earns 0}
\end{rubric}

% ---- FRQ 2 (short) ---------------------------------------------------------
\frq{short}{4}{Zumdahl \S15.2--15.3 \textbullet\ buffer action}

\SI{1.00}{\litre} of a buffer is \SI{0.30}{\Molar} in \ce{HC2H3O2} and
\SI{0.30}{\Molar} in \ce{NaC2H3O2}.

\begin{parts}
  \item Calculate the initial pH. \answerblank[3cm]{4.74}
  \item \SI{0.060}{\mole} of \ce{NaOH} is added. Write the reaction that
        occurs and calculate the new pH. \workspace[2.8cm]
        \keyonly{\begin{apcanswer}\ce{OH^- + HC2H3O2 -> C2H3O2^- + H2O};
        \ce{HC2H3O2} $= 0.24$, \ce{C2H3O2^-} $= 0.36$ mol;
        pH $= 4.74 + \log(0.36/0.24) = 4.74 + 0.18 =
        \mathbf{4.92}$\end{apcanswer}}
  \item The same \SI{0.060}{\mole} of \ce{NaOH} added to
        \SI{1.00}{\litre} of pure water gives pH 12.78. Account for the
        difference. \answerlines[3]{The buffer converts the added
        \ce{OH-} into acetate ion by reaction with acetic acid, so
        $[\ce{OH-}]$ never accumulates and only the ratio of the pair
        changes. In pure water there is nothing to react with the base, so
        all of it remains as free hydroxide.}
\end{parts}

\begin{rubric}
  \rpoint{1}{(a) pH $= 4.74$}
  \rpoint{2}{(b) 1 pt balanced reaction with correct new amounts; 1 pt
    pH $= 4.92$}
  \rpoint{1}{(c) explanation must say the \ce{OH-} is \emph{consumed by the
    weak acid} --- ``buffers resist pH change'' restates the question and
    earns 0}
\end{rubric}

\examtotal

\end{document}
